\documentclass[12pt]{article}

% =========================
% 0) Metadata (EDIT ME)
% =========================
\newcommand{\CourseCode}{MATH 431}
\newcommand{\CourseTitle}{Introduction to Probability}
\newcommand{\CourseSection}{LEC 002} % change to 001 or 003 if needed
\newcommand{\Semester}{Spring 2026}
\newcommand{\Instructor}{Sarah Tammen} % LEC 001: Ander Aguirre | LEC 002: Sarah Tammen | LEC 003: Nicholas Derr
\newcommand{\MeetingInfo}{MWF 9:55--10:45, 6102 Sewell Social Sciences} % LEC 001 default; update if in 002/003
\newcommand{\HomeworkNumber}{12}
\newcommand{\YourName}{Alejandro De La Torre}
\newcommand{\DueDate}{May 1, 2026} % or e.g. February 2, 2026

% =========================
% 1) Core math tools
% =========================
\usepackage{amsmath, amssymb, amsthm}

% =========================
% 2) Lists and spacing
% =========================
\usepackage{enumitem}
\setlist[enumerate,1]{label=\textbf{(\alph*)}, leftmargin=2em, itemsep=0.25em}
\setlist[itemize]{leftmargin=1.5em, itemsep=0.25em}
\setlength{\parskip}{0.35em}
\setlength{\parindent}{0pt}

% =========================
% 3) Page geometry + header/footer
% =========================
\usepackage{geometry}
\geometry{margin=1in}

\usepackage{fancyhdr}
\pagestyle{fancy}
\fancyhf{}
\lhead{\CourseCode\ --- \CourseSection, \Semester}
\rhead{Homework \HomeworkNumber}
\cfoot{\thepage}
\renewcommand{\headrulewidth}{0.4pt}
\setlength{\headheight}{15pt}
\addtolength{\topmargin}{-2pt}

% =========================
% 4) Links (subtle)
% =========================
\usepackage[hidelinks]{hyperref}
\setcounter{tocdepth}{1}

% =========================
% 5) Figures (optional)
% =========================
\usepackage{graphicx}
\graphicspath{{images/}}

% =========================
% 6) Lightweight boxes (optional)
% =========================
\usepackage{mdframed}
\newmdenv[
  skipabove=0.6\baselineskip,
  skipbelow=0.6\baselineskip,
  linewidth=0.6pt,
  roundcorner=4pt,
  innerleftmargin=10pt,
  innerrightmargin=10pt,
  innertopmargin=8pt,
  innerbottommargin=8pt
]{boxnote}

% =========================
% 7) Probability / notation macros
% =========================
\newcommand{\R}{\mathbb{R}}
\newcommand{\N}{\mathbb{N}}
\newcommand{\Z}{\mathbb{Z}}

\newcommand{\OmegaSpace}{\Omega}
\newcommand{\FField}{\mathcal{F}}
% Probability measure in case it is relevant or want to distinguish it from P(A)
%\newcommand{\PMeasure}{\mathbb{P}} 
\newcommand{\E}{\mathbb{E}}

\newcommand{\given}{\,\middle|\,}
\newcommand{\ind}{\mathbf{1}}

% Common “defined as” symbol
\newcommand{\defeq}{\vcentcolon=}

% =========================
% 8) Problem command (TOC + PDF bookmarks)
% Usage:
%   \problem{<label>}
%   \problem[Short title]{<label>}
% Example:
%   \problem{1.4}
%   \problem[Sample space]{1.4}
% =========================
\makeatletter
\newcommand{\problem}[2][]{%
  \def\prob@title{Problem #2\if\relax\detokenize{#1}\relax\else\ (\textit{#1})\fi}%
  \phantomsection%
  \pdfbookmark[1]{\prob@title}{prob#2}%
  \addcontentsline{toc}{section}{\prob@title}%
  \section*{\prob@title}%
  \vspace{-0.25em}\hrule\vspace{0.8em}%
}
\makeatother

% =========================
% 9) Solution environment
% =========================
\newenvironment{solution}{%
  \noindent\textbf{Solution.}\ \ignorespaces
}{\par}

% Optional scratch / planning block
\newenvironment{scratch}{%
  \begin{boxnote}\textbf{Scratch / Setup.}\ \small
}{\end{boxnote}}

% Final answer command: expects math only (no $$ or \[ \])
\newcommand{\finalanswer}[1]{%
  \par\medskip
  \noindent\textbf{Final:}\ \fbox{$#1$}\par
} % AGAIN: MATH ONLY INSIDE THE BOX; if word answer, use \text{} inside or use \fbox{TEXT} or \framebox

% =========================
% 10) Title block
% =========================
\title{Homework \HomeworkNumber}
\author{\YourName \\
\CourseCode\ --- \CourseTitle \\
\CourseSection \\
Instructor: \Instructor}
\date{Due: \DueDate}

\begin{document}
\maketitle

% \section*{Collaboration / Resources}
% (Optional) Write “I worked alone” or list collaborators/resources.

\tableofcontents
\bigskip
\hrule
\bigskip
\newpage

% ============================================================
% Problems (duplicate blocks as needed)
% ============================================================



\problem{9.2}
Let $X$ be an exponential random variable with parameter $\lambda = \tfrac{1}{2}$.
\begin{enumerate}[label=\textbf{(\alph*)}]
\item Use Markov’s inequality to find an upper bound for $P(X > 6)$.
\item Use Chebyshev’s inequality to find an upper bound for $P(X > 6)$.
\item Explicitly compute the probability above and compare with the upper bounds you derived.
\end{enumerate}

\begin{solution}

\begin{scratch}
For $X \sim \mathrm{Exp}(\lambda)$ with $\lambda = \tfrac{1}{2}$,
\[
E[X] = \frac{1}{\lambda} = 2,
\qquad
\operatorname{Var}(X) = \frac{1}{\lambda^2} = 4.
\]
Relevant facts:
\begin{itemize}
\item $X \ge 0$ almost surely, so Markov's inequality applies.
\item Chebyshev's inequality says
\[
P\bigl(|X - E[X]| \ge t\bigr) \le \frac{\operatorname{Var}(X)}{t^2}.
\]
\item Since $E[X] = 2$, the event $\{X > 6\}$ implies $\{|X - 2| \ge 4\}$.
\item For an exponential random variable,
\[
P(X > x) = e^{-\lambda x}, \qquad x \ge 0.
\]
\end{itemize}
\end{scratch}

\begin{enumerate}[label=\textbf{(\alph*)}]
\item Since $X \ge 0$, Markov's inequality gives
\[
P(X > 6) \le \frac{E[X]}{6} = \frac{2}{6} = \frac{1}{3}.
\]
\finalanswer{P(X > 6) \le \frac{1}{3}}

\item Because $\{X > 6\} \subseteq \{|X - 2| \ge 4\}$, we have
\[
P(X > 6) \le P(|X - 2| \ge 4).
\]
Now apply Chebyshev's inequality:
\[
P(|X - 2| \ge 4) \le \frac{\operatorname{Var}(X)}{4^2}
= \frac{4}{16}
= \frac{1}{4}.
\]
Therefore
\[
P(X > 6) \le \frac{1}{4}.
\]
\finalanswer{P(X > 6) \le \frac{1}{4}}

\item The exponential tail can be computed exactly:
\[
P(X > 6) = e^{-\lambda \cdot 6} = e^{-3} \approx 0.0498.
\]
Comparing the three values,
\[
e^{-3} \approx 0.0498 < \frac{1}{4} < \frac{1}{3}.
\]
So both inequalities give valid upper bounds, and the Chebyshev bound is better than the Markov bound because $\frac{1}{4}$ is smaller and closer to the exact probability.
\finalanswer{P(X > 6) = e^{-3} \approx 0.0498,\quad \text{and Chebyshev gives the better upper bound}}
\end{enumerate}

\end{solution}

\newpage

\problem{9.6}
Nate is a competitive eater specializing in eating hot dogs. From past experience we know that it takes him on average $15$ seconds to consume one hot dog, with a standard deviation of $4$ seconds. In this year’s hot dog eating contest he hopes to consume $64$ hot dogs in just $15$ minutes. Use the CLT to approximate the probability that he achieves this feat of skill.

\begin{solution}

\begin{scratch}
Let $X_i$ be the time, in seconds, that Nate needs to consume the $i$th hot dog, and let
\[
S_{64} = X_1 + \cdots + X_{64}.
\]
Then Nate succeeds exactly when $S_{64} \le 15$ minutes $= 900$ seconds.

From the data,
\[
E[X_i] = 15,
\qquad
\operatorname{SD}(X_i) = 4,
\qquad
\operatorname{Var}(X_i) = 16.
\]
Assuming the $X_i$ are i.i.d.,
\[
E[S_{64}] = 64 \cdot 15 = 960,
\qquad
\operatorname{Var}(S_{64}) = 64 \cdot 16 = 1024,
\qquad
\operatorname{SD}(S_{64}) = 32.
\]
Use the CLT because the exact distribution of $S_{64}$ is not given:
\[
\frac{S_{64} - 960}{32} \approx N(0,1).
\]
\end{scratch}

We want the probability that Nate consumes $64$ hot dogs within $15$ minutes, so
\[
P(\text{Nate succeeds}) = P(S_{64} \le 900).
\]
By the central limit theorem,
\[
P(S_{64} \le 900)
\approx
P\!\left(
\frac{S_{64} - 960}{32}
\le
\frac{900 - 960}{32}
\right)
= \Phi\!\left(\frac{-60}{32}\right)
= \Phi(-1.875).
\]
Therefore,
\[
P(S_{64} \le 900) \approx \Phi(-1.875) \approx 0.0304.
\]
So the CLT approximation says that Nate has about a $3.04\%$ chance of finishing $64$ hot dogs in $15$ minutes.
\finalanswer{P(S_{64} \le 900) \approx \Phi(-1.875) \approx 0.0304}

\end{solution}

\newpage

\problem{9.18}
Suppose that $X \sim \text{NegBin}(100, \tfrac{1}{3})$.
\begin{enumerate}[label=\textbf{(\alph*)}]
\item Give an upper bound for $P(X > 500)$ using Markov’s inequality.
\item Give an upper bound for $P(X > 500)$ using Chebyshev’s inequality.
\item Recall that $X$ can be written as the sum of $100$ i.i.d. random variables with Geom$(\tfrac{1}{3})$ distribution. Use this representation and the CLT to estimate $P(X > 500)$.
\item Recall that $X$ can be represented as the number of trials needed to get the $100$th success in a sequence of i.i.d. trials of success probability $\tfrac{1}{3}$. Estimate $P(X > 500)$ by looking at the distribution of the number of successes within the first $500$ trials.
\end{enumerate}

\begin{solution}

\begin{scratch}
Here $X \sim \mathrm{NegBin}(100,\tfrac{1}{3})$ means that $X$ is the number of trials needed to obtain the $100$th success when each trial succeeds with probability
\[
p = \frac{1}{3}.
\]
With $r = 100$ and this convention,
\[
E[X] = \frac{r}{p} = \frac{100}{1/3} = 300,
\]
\[
\operatorname{Var}(X) = \frac{r(1-p)}{p^2}
= \frac{100(2/3)}{(1/3)^2}
= 600,
\qquad
\operatorname{SD}(X) = \sqrt{600} = 10\sqrt{6}.
\]

For part (c), write
\[
X = G_1 + \cdots + G_{100},
\]
where the $G_i$ are i.i.d.\ $\mathrm{Geom}(\tfrac{1}{3})$ random variables.

For part (d), let $N$ be the number of successes in the first $500$ trials. Then
\[
N \sim \mathrm{Bin}\!\left(500,\frac{1}{3}\right),
\qquad
\{X > 500\} = \{N \le 99\}.
\]
Also,
\[
E[N] = \frac{500}{3},
\qquad
\operatorname{Var}(N) = 500 \cdot \frac{1}{3} \cdot \frac{2}{3} = \frac{1000}{9},
\qquad
\operatorname{SD}(N) = \sqrt{\frac{1000}{9}}.
\]
\end{scratch}

\begin{enumerate}[label=\textbf{(\alph*)}]
\item Since $X \ge 0$, Markov's inequality applies:
\[
P(X > 500) \le \frac{E[X]}{500} = \frac{300}{500} = \frac{3}{5}.
\]
\finalanswer{P(X > 500) \le \frac{3}{5}}

\item Since $X > 500$ implies $X - 300 > 200$, we have
\[
\{X > 500\} \subseteq \{|X - 300| \ge 200\}.
\]
Hence Chebyshev's inequality gives
\[
P(X > 500)
\le
P(|X - 300| \ge 200)
\le
\frac{\operatorname{Var}(X)}{200^2}
= \frac{600}{40000}
= \frac{3}{200}
= 0.015.
\]
\finalanswer{P(X > 500) \le \frac{3}{200} = 0.015}

\item Using the representation $X = G_1 + \cdots + G_{100}$ and the CLT,
\[
\frac{X - 300}{\sqrt{600}} \approx N(0,1).
\]
Therefore
\[
P(X > 500)
=
P\!\left(\frac{X - 300}{\sqrt{600}} > \frac{500 - 300}{\sqrt{600}}\right)
\approx
1 - \Phi\!\left(\frac{200}{\sqrt{600}}\right).
\]
Since
\[
\frac{200}{\sqrt{600}} = \frac{20}{\sqrt{6}} \approx 8.165,
\]
we get
\[
P(X > 500) \approx 1 - \Phi\!\left(\frac{20}{\sqrt{6}}\right) \approx 1.6 \times 10^{-16}.
\]
\finalanswer{P(X > 500) \approx 1 - \Phi\!\left(\frac{20}{\sqrt{6}}\right) \approx 1.6 \times 10^{-16}}

\item Let $N$ be the number of successes in the first $500$ trials. Then
\[
N \sim \mathrm{Bin}\!\left(500,\frac{1}{3}\right)
\]
and
\[
P(X > 500) = P(N \le 99).
\]
Using the CLT with a continuity correction,
\[
P(N \le 99)
\approx
P(N < 99.5)
=
P\!\left(
\frac{N - 500/3}{\sqrt{1000/9}}
<
\frac{99.5 - 500/3}{\sqrt{1000/9}}
\right)
\approx
\Phi\!\left(\frac{99.5 - 500/3}{\sqrt{1000/9}}\right).
\]
Now
\[
\frac{99.5 - 500/3}{\sqrt{1000/9}} \approx -6.372,
\]
so
\[
P(X > 500) = P(N \le 99) \approx \Phi(-6.372) \approx 9.3 \times 10^{-11}.
\]
This is different from part (c), which is not surprising: both are CLT approximations, and this is an extreme tail event.
\finalanswer{P(X > 500) = P(N \le 99) \approx \Phi(-6.372) \approx 9.3 \times 10^{-11}}
\end{enumerate}

\end{solution}


\end{document}
